\section{Introduction and problem setting}\label{sec:foundations}

Many leader--follower models contain many local follower choices but only
a few shared resources, aggregate quantities and leader decisions. A tariff
may influence the operation of many units whose nonlinear interaction
depends only on total load or emissions; a local variable or small block
describes the choices of one unit. We model one follower that
optimizes the joint allocation. This also covers separate agents whose
decisions decouple, but a nonconvex aggregate term does not in general
describe a Nash equilibrium among independent agents.

We ask whether the number of shared quantities, rather than the number of
follower variables, can control the complexity of \emph{global} bilevel
optimization. Linear bilevel optimization is strongly
NP-hard~\cite{HansenJaumardSavard1992}; see also~\cite{Buchheim2023,KetkovProkopyev2026}.
A linear class remains NP-complete with a single leader
variable~\cite{SugishitaCarvalhoIPCO2026,SugishitaCarvalho2026}. Positive
results fix the number of follower
variables~\cite{LiuSpencer1995,Deng1998,KetkovProkopyev2026} or, for linear
followers, the number of follower constraints~\cite{KetkovProkopyev2026}. A tractable follower
problem does not suffice: the leader optimizes over the responses induced
by all its decisions, and this problem can be nonconvex even when every
follower has a unique convex minimizer. Standard single-level
reformulations do not answer the question either. A KKT reformulation
replaces follower optimality by stationarity. This is exact for convex
followers under a constraint qualification, which linear constraints such
as ours satisfy, but for nonconvex followers stationary points need not be
global optima; its mixed-integer form also needs valid big-$M$
bounds~\cite{KleinertEtAl2021,Buchheim2023}.
A value-function reformulation is exact for nonconvex followers but
quantifies over the entire follower vector~\cite{JeyakumarEtAl2016,NieWangYe2017}.
We therefore separate three tasks: describing all global follower
responses, optimizing over them under a stated optimistic or pessimistic
selection convention, and representing the output at a stated accuracy.

\paragraph{Key idea.}
In the model of \cref{subsec:model}, fix a leader $x$ and an aggregate
value $w=U(x)z$. On the fiber $\{z\in P(x):U(x)z=w\}$ the aggregate cost is
constant and the remaining local cost is strictly convex, so a global
follower optimum is the unique minimizer on its own fiber. With multipliers
for the $k$ shared rows, each block becomes a small strictly convex
parametric quadratic program whose solution is piecewise rational in the
compressed coordinates $\zeta=(x,w,\lambda,\mu)\in\R^{r+s+k}$, with
polynomially many branches per block. We call
$r+s+k$ the \emph{response description dimension}. Only polynomially many
joint regimes of all blocks are realizable in this fixed dimension, and a
value-function comparison in the compressed coordinates, rather than in all
$N$ follower variables, selects exactly the global responses.

\subsection{Contributions and scope}\label{subsec:contributions}

Under explicit local structure, a large follower dimension can thus be
replaced by a fixed response description dimension. Our original
contributions are the combined complexity and recovery guarantees
below, together with the restriction-preserving boundary constructions.
Degrees are numerical degrees, and exponents may depend on the fixed
dimensions (\cref{subsec:encoding}).

\paragraph{Exact global responses (\cref{sec:exact}).}
The follower has positive-definite quadratic local blocks of fixed
dimension $d$, $k$ shared linear rows, and a polynomial, possibly nonconvex
cost of $s$ aggregates; the leader has $r$ variables. The numbers of
blocks, local rows and upper polynomial rows, and the supports of the upper
polynomials, can grow. For fixed $(r,d,s,k)$, \cref{thm:exact-block}
decides optimistic feasibility and returns an attained optimizer and its
value in one real-algebraic field of polynomial encoding length, in time
polynomial in the input length and numerical degree. \Cref{thm:exact-semantics}
extends this to polynomially moving constraint normals and to pessimistic
selection with universal upper feasibility, computing exact infima and
deciding attainment. Fixed normals give optimistic attainment
(\cref{lem:closed-response}); moving normals or pessimistic selection can
lose it (\cref{ex:moving-nonattainment,ex:pessimistic-nonattainment}).
With constant local Hessians and fixed input degree, the algebraic degree
of the output is independent of the number of follower variables, for
fixed normals (\cref{cor:constant-hessian-degree}) and, in the robust
setting, with constant local rows but moving shared and measurement normals
(\cref{cor:robust-degree}). For a
supplied diagonal-plus-fixed-rank Hessian and affine upper data, rational
linear programs suffice (\cref{cor:low-rank-lp}). \Cref{thm:fixed-core}
gives a related polyhedral elimination theorem, with recovery of all blocks
by one common convex combination of a few support tuples.

\paragraph{Near-optimality robustness (\cref{sec:robust-screening}).}
Cost-near-optimal followers need not be stationary, so we minimize on the
fibers of the affine measurements used by each upper criterion.
\Cref{thm:nearoptimal-robust} gives exact robust feasibility and
worst-objective optimization when every criterion uses at most a fixed
number $h$ of measurements, even if the combined measurement rank grows.
Without this bound, computing the worst value of one convex quadratic
criterion is NP-hard (\cref{prop:robust-measurement-hardness}). Attainment
holds for convex followers with fixed normals (\cref{prop:robust-attainment})
but can fail for a nonconvex follower with a unique nominal optimum and a
positive budget (\cref{ex:positive-budget-nonattainment}).

\paragraph{Certified recovery of a dense quadratic follower (\cref{sec:robust-screening}).}
A perturbation argument certifies coordinate statuses of a dense convex
box-constrained quadratic response over whole cells of a structured
surrogate's response cover. If the cover has $M$ cells and at most $t$
uncertified coordinates per cell, exact bilevel recovery needs at most
$M3^t$ LPs after screening (\cref{thm:screening-recovery}). A computable
neighborhood of the surrogate bounds $t$ under a transition-multiplicity
condition (\cref{cor:screening-neighborhood}); a small matrix error alone
is insufficient (\cref{ex:screening-conservatism}).

\paragraph{Global approximation in accuracy bits (\cref{sec:accuracy}).}
For strictly convex polynomial local costs, affine resource data and a
convex polynomial aggregate, \cref{thm:accuracy-global} returns a rational
leader with global additive error $2^{-B}$ for an arbitrary explicitly
encoded polynomial upper objective, in time polynomial in input length,
numerical degree and $B$ for fixed $r,k,s$. No uniform positive curvature
or Slater condition is assumed. The proof combines uniform rational inverse
approximations (\cref{thm:inverse-approximation}), residual-to-response
bounds, and rational recovery that preserves the base feasible leader
polytope. Separately, the response is H\"older continuous in the leader
with the sharp exponent $1/P$, where $P$ is the largest local marginal
degree (\cref{prop:accuracy-response-modulus}). Response-dependent upper
inequalities admit outer and inner guarantees (\cref{thm:accuracy-upper}).
Exactly feasible near-optimal rational output requires a tightening
modulus, a convex strict anchor or reserve controls
(\cref{cor:accuracy-tightening,cor:accuracy-anchor,cor:accuracy-reserve});
\cref{ex:accuracy-isolated-optimum,ex:accuracy-irrational-feasibility}
show that some such condition is necessary.

\paragraph{Structural and arithmetic boundaries (\cref{sec:boundaries}).}
Exact optimization is NP-hard already with one leader, a unit-box
follower, an affine upper objective and a dense positive-definite Hessian
with condition number below two and arbitrarily close to the identity
(\cref{thm:conditioned-hardness}). A construction with a constant gap
(\cref{thm:dense-hardness}) rules out, unless $\mathrm P=\mathrm{NP}$,
exactly feasible solutions with multiplicative ratio below three
(\cref{cor:dense-multiplicative}). For fixed leader dimension and no
response-dependent upper rows, a complementary scheme runs in time
polynomial in the input length, the condition measure
$\norm Q_\infty\norm{Q^{-1}}_\infty$ and the inverse normalized additive
error, not its bit length; this is consistent with the exact hardness
(\cref{thm:conditioned-grid}). For growing leader dimension we transfer
W[1]-hardness and an ETH lower bound (\cref{thm:growing-leader}), give a
polynomial algorithm for a supplied bounded leader core with a diagonal
follower Hessian and no response-dependent upper rows
(\cref{thm:leader-core}), and show weak NP-hardness when the leader
interaction graph is a path (\cref{thm:leader-path}). A path in the
follower \emph{constraints} gives NP-completeness with one leader, an
identity Hessian and a fixed alphabet of local constraint coefficients;
this reduction retains one explicitly nonconvex upper quadratic row
(\cref{thm:follower-path-hardness}). Without positive local curvature,
optimistic feasibility is NP-complete even with no leader and independent
unit boxes, using quadratic upper rows with zero local curvature or linear
upper rows with negative local curvature (\cref{prop:curvature-boundary}).
Exponential algebraic degree (\cref{prop:root-sum-degree}) and
$\Omega(P)$-bit rational-output lower bounds (\cref{prop:single-power-output})
show why exact field output, numerical degree and accuracy bits must be
treated separately.

\paragraph{Complete scalar computations (\cref{sec:computation}).}
We implement an exact certified convex path method
(\cref{prop:convex-certified-path}) and an exact nonconvex scalar atlas
that reconstructs original responses and solves both upper selection
conventions, including attainment (\cref{thm:nonconvex-atlas}).
Convexification alone admits false follower choices
(\cref{ex:false-contacts}); the atlas retains original contacts. An
independent original-coordinate baseline enumerates stationary faces in original coordinates
and solves the same continuous leader task under a checked nonsingularity
restriction (\cref{prop:original-face-baseline}); it is a validation
baseline, not a new face-enumeration principle. On the tested small
instances both solvers agree on every exact value and attainment flag, and
neither is uniformly faster. The experiments separate heterogeneous
nonconvex instances from large repeated populations and report a negative
result: certified dense screening, including its preprocessing, was slower
than a numerical MILP comparator on every tested instance. The
implementations cover these specializations only, not the general
real-algebraic algorithms.

\subsection{Related work}\label{subsec:related}

\paragraph{Complexity with fixed dimensions.}
Liu and Spencer~\cite{LiuSpencer1995} give a polynomial algorithm for
optimistic linear bilevel optimization with a fixed number of
\emph{follower} variables, and Deng~\cite{Deng1998} discusses the resulting
complexity theory. Fixing the number of leader variables while the follower
dimension grows is different. Sugishita and
Carvalho~\cite{SugishitaCarvalhoIPCO2026,SugishitaCarvalho2026} prove
NP-completeness with one leader variable, unit-bounded variables and no
additional upper constraints for a linear bilevel class with coupled
follower constraints; their polynomial algorithm for a local optimum does
not give global tractability. Ketkov and Prokopyev~\cite{KetkovProkopyev2026}
distinguish fixed follower variable and constraint counts, optimistic and
pessimistic responses, and linear and quadratic models. Their optimistic
convex-quadratic positive result fixes the follower dimension and imposes
their stated convex quadratic upper structure. Neither that result nor a
bound on the total follower constraint count covers arbitrarily many local
follower variables and local bounds with few shared constraints and a
possibly nonconvex polynomial aggregate cost; counting shared rows
separately from local rows is essential. Buchheim~\cite{Buchheim2023} proves
NP membership for linear bilevel optimization; polynomial-size certificates
need not give polynomial-time optimization.

\paragraph{Optimistic and pessimistic selection.}
The pessimistic convention of~\cite{KetkovProkopyev2026} illustrates why
upper feasibility must be specified before complexity statements are
compared. Henke, Lefebvre, Schmidt and Th\"urauf~\cite{HenkeEtAl2026} use
nonempty responses and universal upper-row feasibility for pessimistic
linear bilevel optimization, as we do in \eqref{eq:pessimistic-domain}.
Under their standing assumptions, they give polynomial-size reformulations
between linear variants that preserve projected global optimizers. Their
Theorem~3.3 introduces a full follower-sized vector among the leader
variables, so it does not preserve a fixed leader dimension as the follower
dimension grows and does not directly give the bounds proved here.

\paragraph{KKT and value-function methods for polynomial bilevel problems.}
Jeyakumar et al.~\cite{JeyakumarEtAl2016} and Nie, Wang and
Ye~\cite{NieWangYe2017} develop convergent semidefinite relaxations under
their respective assumptions; the latter use globally quantified
follower-value comparisons in a semi-infinite reformulation. Nie et
al.~\cite{NieEtAl2021} represent multipliers by polynomial or rational
expressions and combine them with feasible extensions and exchange steps.
Nie, Ye and Zhong~\cite{NieYeZhong2026} use sparse multiplier supports and
disjunctive decompositions for linearly constrained followers. Thus global
value comparison, rational multiplier expressions and conic support
reduction are not new here. The difference is quantitative: we also
eliminate the \emph{primal} local variables and bound the number of joint
regimes in a dimension independent of their total count. A support bound
based on the rank of the entire follower constraint matrix can grow like
the follower dimension when that matrix contains all box normals; our local
support bound uses the fixed block dimension. Convergence of global
polynomial subproblems and a polynomial rational-bit bound are different
conclusions.

\paragraph{Separable and parametric quadratic programming.}
Megiddo and Tamir~\cite{MegiddoTamir1993} use low-dimensional multiplier
arrangements for separable convex quadratic optimization and for
fixed-dimensional quadratic blocks with shared constraints; our local
formulas and multiplier partitions build on this principle. The block
branches follow the active-set mechanism of explicit parametric quadratic
programming~\cite{BemporadEtAl2002}, here with polynomial parameter
dependence. Exact sign determination and quantifier elimination in fixed
dimension are classical~\cite{BasuPollackRoy1996}. Fixed rank of the entire
quadratic matrix is another tractable restriction for box-constrained
optimization~\cite{HladikCernyRada2021}. It differs from a positive
diagonal or block-diagonal matrix plus fixed-rank coupling, whose total
rank may grow with the follower count; our result uses the local curvature
and the constraint structure, not low-rank nonconvexity alone.

\paragraph{Approximation.}
Hochbaum and Shanthikumar~\cite{HochbaumShanthikumar1990} give logarithmic
accuracy dependence for approximating the solution vector of continuous
separable allocation under their function-oracle assumptions, with
numerical dependence on a bound for constraint subdeterminants.
Vigneron~\cite{Vigneron2014} gives fully polynomial approximation schemes
for global optimization of sums of nonnegative algebraic functions of
constant description complexity, with polynomial dependence on inverse
relative error. A high-accuracy follower evaluation at a fixed leader
differs from a global approximation algorithm for the induced, possibly
signed upper objective; the comparison must track both numerical degree and
accuracy dependence. Among first-order methods, Chen, Ji and
Zhang~\cite{ChenJiZhang2026} study condition-number dependence for
approximate stationarity with a nonconvex upper problem and a strongly
convex lower problem. Wu et al.~\cite{WuEtAl2026} study stationarity with a
uniformly convex unconstrained follower; their Lemma~4.2 already gives the
response exponent $1/(p-1)$ for uniform convexity of order $p$. Our $1/P$
bound extends this exponent to the structured polynomial model with moving
affine resources and a constant of polynomial encoding length. Xiao and
Chen~\cite{XiaoChen2025} prove global convergence of penalty-based bilevel
gradient descent, with approximate lower optimality, under joint or
blockwise Polyak--\L{}ojasiewicz conditions on the penalty reformulation
and further assumptions; our guarantees use fixed structural dimensions instead of
these landscape conditions. Certified approximate parametric paths and
optimizer-error bounds~\cite{GiesenJaggiLaue2012,GiesenEtAl2012,BemporadFilippi2001}
are earlier treatments related to \cref{thm:conditioned-grid}.

\paragraph{Near-optimality robustness and screening.}
Besan\c{c}on, Anjos and Brotcorne~\cite{BesanconAnjosBrotcorne2024,BesanconAnjosBrotcorne2021}
study robustness of upper feasibility against cost-near-optimal responses
and its complexity; Beck, Ljubi\'c and Schmidt~\cite{BeckLjubicSchmidt2023}
place this model in the broader uncertainty literature. Our model also
allows a worst-case upper objective; the contribution is exact compression
of its adversarial problems, not the near-optimality model itself.
Surrogate-based safe screening has a direct
precedent~\cite{DantasGribonval2019}, and variational-inequality screening
enclosures precede it~\cite{LiuZhaoWangYe2014}; our target is an exact
global upper optimum for the true dense follower.

\paragraph{Boundary sources and large populations.}
Several boundary constructions in \cref{sec:boundaries} adapt known
ones~\cite{SugishitaCarvalho2026,MairalYu2012,FroeseEtAl2026,GartnerEtAl2013,PardalosVavasis1991},
as attributed there; in particular, the W[1]/ETH result for growing leader
dimension is transferred from~\cite{FroeseEtAl2026}. Flocco, Schiewe and
Gabriel~\cite{FloccoEtAl2026} develop nested Benders decomposition for many
linear followers that decouple at a fixed leader; their model and
computational objective differ from ours.

Table~\ref{tab:prior-comparison} summarizes the closest comparisons. To our
knowledge, the cited work does not establish the combined exact
quadratic-block or global accuracy-bit guarantees with unbounded follower
dimension and fixed shared dimensions; we claim no priority for the
ingredients (multiplier arrangements, global value comparison, polynomial
inversion, convex conjugation, fixed-dimensional real algebra).

\begin{table}[t]
\centering\small
\caption{Relationship to selected prior work; the right column gives the
additional guarantee under this paper's restrictions.}
\label{tab:prior-comparison}
\begin{tabular}{@{}>{\raggedright\arraybackslash}p{.23\linewidth}>{\raggedright\arraybackslash}p{.32\linewidth}>{\raggedright\arraybackslash}p{.39\linewidth}@{}}
\toprule
Prior work & Established result or method & Distinction in this paper\\\midrule
Liu--Spencer; Ketkov--Prokopyev
\cite{LiuSpencer1995,KetkovProkopyev2026} &
Positive bilevel results with fixed follower dimension; further distinctions
by total constraint count and selection convention. &
Growing follower dimension and local row counts; fixed leader, local block,
shared resource and aggregate dimensions.\\[5pt]
Megiddo--Tamir; Nie--Ye--Zhong
\cite{MegiddoTamir1993,NieYeZhong2026} &
Low-dimensional quadratic multiplier arrangements; sparse multiplier
supports and KKT disjunctions. &
Polynomially many jointly realizable local regimes; primal recovery and
comparison of all global responses in fixed dimension.\\[5pt]
Hochbaum--Shanthikumar; Vigneron
\cite{HochbaumShanthikumar1990,Vigneron2014} &
High-accuracy separable allocation; approximation schemes for sums of
nonnegative algebraic functions. &
Global signed polynomial upper optimization in accuracy bits, fixed shared
coupling, and rational leader recovery with explicit feasibility guarantees.\\[5pt]
Besan\c{c}on--Anjos--Brotcorne
\cite{BesanconAnjosBrotcorne2024,BesanconAnjosBrotcorne2021} &
Cost-near-optimal response robustness and its complexity framework. &
Exact nonconvex measurement-fiber compression, with a fixed measurement count
separately for each criterion.\\[5pt]
Giesen et al.; Bemporad--Filippi
\cite{GiesenJaggiLaue2012,GiesenEtAl2012,BemporadFilippi2001} &
Certified approximate parametric paths and optimizer-error bounds. &
A saturation-based cost grid with a condition-dependent global upper bound
and no numerical incentive-magnitude factor beyond input encoding.\\[5pt]
Gardiner--Lucet; Moehle et al.
\cite{GardinerLucet2010,MoehleEtAl2023} &
Univariate piecewise quadratic convex-envelope constructions. &
An exact scalar bilevel specialization retaining original contacts, upper
feasibility, selection and attainment, with reproducible comparisons.\\
\bottomrule
\end{tabular}
\end{table}

\paragraph{Organization and notation.}
Sections~\ref{sec:exact}--\ref{sec:computation} follow the order of the
contributions above, and Section~\ref{sec:conclusions} concludes. The
appendices give full proofs for polyhedral elimination, inverse
approximation, quantitative resource bounds and path geometry. The principal
dimensions are $r$ (leader), $N$ (total follower), $d$ (maximum local
block), $s$ (aggregate) and $k$ (shared rows). Only the parameters
explicitly fixed in each theorem are held constant. Matrix symbols are local
to their model; in particular $H$ denotes shared inequality normals in the
quadratic-block model and an upper polynomial in Section~\ref{sec:accuracy}.

\subsection{The quadratic-block model}\label{subsec:model}

All polynomial data are supplied explicitly with rational coefficients.
The leader chooses $x\in C\subset\R^r$, where $C$ is a compact set
specified by a rational quantifier-free polynomial formula. The follower
vector is $z=(z_1,\ldots,z_B)$, with $z_b\in\R^{d_b}$,
$d_b\le d$, and total dimension $N=\sum_{b=1}^B d_b$. Both $B$ and $N$
may grow. For the primary model, let
\begin{align}
 P_b(x)&=\{y\in\R^{d_b}:E_b y=e_b(x),\ G_b y\le h_b(x)\},
 \label{eq:local-polytope}\\
 P(x)&=\left\{z:z_b\in P_b(x)\ (b=1,\ldots,B),\
                     Az=b(x),\ Hz\le a(x)\right\}.
 \label{eq:follower-polytope}
\end{align}
The matrices $E_b,G_b,A,H$ are rational and independent of $x$.
There are $k_{=}$ shared equality rows and $k_{\le}$ shared inequality
rows, and $k=k_=+k_{\le}$. The number of local rows is unrestricted.
All right-hand sides are polynomial in $x$. Among the local inequalities
are coordinate bounds
$\ell_{bi}(x)\le z_{bi}\le u_{bi}(x)$, with polynomial endpoints
satisfying $\ell_{bi}(x)\le u_{bi}(x)$ on $C$.
Hence the union of the follower feasible sets is bounded.
A set $P_b(x)$ or $P(x)$ may be empty; follower feasibility is not
assumed at every leader.

Let $U(x)\in\Q[x]^{s\times N}$, and partition its columns into blocks
$U_b(x)$. For rational polynomial vectors $c_b(x)\in\Q[x]^{d_b}$,
the follower objective is
\begin{equation}\label{eq:quadratic-objective}
 f(x,z)=\sum_{b=1}^B\left(\frac12 z_b\trans Q_b(x)z_b+
                c_b(x)\trans z_b\right)+\phi(x,U(x)z).
\end{equation}
Each $Q_b(x)$ is a symmetric polynomial matrix positive definite for
every $x\in C$. The polynomial $\phi(x,w)$ can be nonconvex in $w$.
Its argument $w=U(x)z\in\R^s$ collects the shared aggregate quantities.
Positive definiteness is imposed on the local blocks, not on the entire
follower Hessian. In particular, the follower need not have a unique
minimizer.

The upper objective $F(x,z)$ and upper rows $g_j(x,z)\le0$
($j=1,\ldots,m$) are explicitly encoded rational polynomials.
Their monomials may involve many different follower coordinates; neither
$m$ nor their support sizes are structural parameters. Unless a theorem
says otherwise, the fixed parameters of the exact quadratic model are
\begin{equation}\label{eq:structural-parameters}
                    (r,d,s,k_=,k_{\le}).
\end{equation}
The scalar local model has $d=1$ and only interval constraints within a
block. Polynomially varying constraint normals are treated in
\cref{subsec:moving-normals}. The approximation results of
\cref{sec:accuracy} use a separate model with strictly convex univariate
polynomial local costs, a convex aggregate cost, and their own
leader-domain and data assumptions; it is not an implicit extension of
\eqref{eq:quadratic-objective}.

\subsection{Follower selection and upper feasibility}\label{subsec:semantics}

For $P(x)\ne\varnothing$, compactness gives the value and solution set
\begin{equation}\label{eq:follower-value}
 v(x)=\min_{z\in P(x)} f(x,z),\qquad
 S(x)=\{z\in P(x):f(x,z)=v(x)\}.
\end{equation}
Put $S(x)=\varnothing$ if $P(x)=\varnothing$, and use $v(x)$ only
on feasible leaders. The follower always optimizes over $P(x)$ without
upper constraints. Define the optimistic feasible set and value by
\begin{align}
 \mathcal O&=\{(x,z):x\in C,\ z\in S(x),\
                             g_j(x,z)\le0\ (j=1,\ldots,m)\},
 \label{eq:optimistic-set}\\
 \operatorname{val}_{\rm opt}&=\inf_{(x,z)\in\mathcal O} F(x,z).
 \label{eq:optimistic-value}
\end{align}
Thus the leader may select a favorable global response that satisfies its
upper rows. A minimizer is asserted only when attainment is established.

For pessimistic optimization, upper feasibility is universal:
\begin{align}
 D_0&=\{x\in C:S(x)\ne\varnothing,\
             g_j(x,z)\le0\text{ for all }z\in S(x),\ j=1,\ldots,m\},
 \label{eq:pessimistic-domain}\\
 W_0(x)&=\max_{z\in S(x)}F(x,z),\qquad
 \operatorname{val}_{\rm pes}=\inf_{x\in D_0}W_0(x).
 \label{eq:pessimistic-value}
\end{align}
The maximum at a fixed feasible leader is attained; the infimum over
leaders need not be. An upper constraint cannot discard an inconvenient
follower optimum before $W_0$ is computed.

Given a nonnegative polynomial budget $\rho(x)$ on $C$, the set of
cost-near-optimal responses is
\begin{equation}\label{eq:nearoptimal-set}
 S_\rho(x)=\{z\in P(x):f(x,z)\le v(x)+\rho(x)\}
\end{equation}
at feasible leaders, and is empty otherwise. Replacing $S(x)$ by
$S_\rho(x)$ in \eqref{eq:pessimistic-domain}--\eqref{eq:pessimistic-value}
defines $D_\rho$, $W_\rho$, and $\operatorname{val}_\rho$.

For revenue maximization, the optimistic objective is a supremum and the
pessimistic objective maximizes the minimum revenue over all responses;
equivalently, negate revenue in the conventions above. The exact
algorithms report infeasibility separately and state whether the finite
optimal value is attained.

\subsection{Input length, degree, and output}\label{subsec:encoding}

The bit model counts arithmetic on binary-encoded rational data. Let $L$
be the total rational input encoding length, including dimensions, row
indices, monomial lists, exponents, and the Boolean formula for $C$.
Let $\delta\ge2$ bound the actual total degree of every input polynomial.
The exact algorithms have running time polynomial in $(L,\delta)$ for
fixed \eqref{eq:structural-parameters}. Polynomial time in $L$ alone
requires the \emph{degree-encoding condition} that $\delta$ is polynomially
bounded by $L$. It holds for dense or unary-exponent encodings, but not for
sparse monomials with unrestricted binary exponents. Polynomial local costs
in the approximation results are supplied densely; a power-specific result
states its dependence on the numerical power.

A common-field algebraic output consists of a square-free polynomial
$p\in\Z[T]$, a rational isolating interval selecting a real root
$\alpha$, and rational polynomials $p_i,q_i$ such that each returned
coordinate equals $p_i(\alpha)/q_i(\alpha)$ with $q_i(\alpha)\ne0$.
The degree of $p$, all coefficient bit lengths, and the combined encoding
of all coordinates are bounded polynomially in $(L,\delta)$ for the stated
fixed dimensions, hence polynomially in $L$ under the degree-encoding
condition. The polynomial $p$ need not be minimal; it suffices that every
recovered coordinate lies in the real field generated by its chosen root.

An accuracy-bit guarantee takes a positive rational tolerance $\eps$,
whose encoding contributes to the running time; for $\eps=2^{-\beta}$ the
bound is polynomial in $\beta$, not in $1/\eps$. In the base
approximation model, the output is a rational leader that is exactly
feasible for the base leader domain and admits a feasible follower, with
an additive guarantee for the upper objective at the true response, not
at an approximate follower vector. Exact satisfaction of
response-dependent upper rows requires additional conditions; without
them, a theorem must specify its allowed upper violation.

The exponents of the polynomial bounds depend on the structural
parameters. For the decision versions this gives XP-type bounds
$L^{h(r,d,s,k)}$ under the degree-encoding condition. No fixed-parameter
bound $h(r,d,s,k)L^c$ with an absolute exponent $c$ is claimed.

\subsection{Algebraic and polyhedral tools}\label{subsec:tools}

We use the following standard consequences of fixed-dimensional real
algebraic computation~\cite[Theorem 1.3.1 and Section 3.1.3]{BasuPollackRoy1996}.
For a rational polynomial family in a fixed total number of real
variables, realizable sign conditions, including zero signs, can be
enumerated in time polynomial in the number of polynomials, their degrees,
and coefficient bit lengths. A first-order formula with a fixed total
number of free and bound variables can be replaced by an equivalent
quantifier-free formula within the same kind of bound. Nonempty
semialgebraic sets admit algebraic samples of polynomial degree and
encoding length. Applied simultaneously to all free coordinates, sampling
gives the common-field representation described above. The total number
of variable occurrences introduced by quantified copies must be controlled;
fixing only the number of free variables is insufficient.

\begin{samepage}
\begin{lemma}[Substitution in fixed dimension]\label{lem:substitution}
Let $v\in\R^q$, with $q$ fixed, and suppose that $Z_1(v),\ldots,Z_N(v)$
and $D(v)$ have explicit polynomial encodings of polynomial size and degree.
On a region where $D>0$, put $z_i=Z_i/D$. Let $p(x,z)$ be an explicitly
listed rational polynomial of total degree at most $\delta$, with the
coordinates of $x$ among those of $v$. Then
\[
             D(v)^\delta p(x,Z(v)/D(v))
\]
is a polynomial that can be constructed in time polynomial in the input
encoding lengths, the degrees of $Z_i,D$, and $\delta$. Its sign agrees
with that of $p(x,z)$ on the region.
\end{lemma}
\end{samepage}
\begin{proof}
An input monomial $a x^\nu z^\gamma$ becomes
$a x^\nu\prod_iZ_i^{\gamma_i}D^{\delta-|\gamma|}$.
Its degree is at most $\delta(1+T)$ if the degrees of $Z_i,D$ are at
most $T$. A polynomial of degree at most $\delta(1+T)$ in $q$ variables
has at most $\binom{\delta(1+T)+q}{q}$ monomials. This number is
polynomial for fixed $q$. Expanding and adding the explicit input
monomials therefore takes polynomially many rational operations;
coefficient bit lengths grow at most polynomially under these products
and sums. Positivity of $D^\delta$ proves sign preservation.
\end{proof}

For polyhedral limits we use Hoffman's bound~\cite{Hoffman1952}.
Given fixed matrices $M,T$, there is a constant $H_{M,T}$ such that,
whenever $K(a,b)=\{y:My\le a,\ Ty=b\}$ is nonempty,
\begin{equation}\label{eq:hoffman}
 \dist(y,K(a,b))\le H_{M,T}
          \bigl(\norm{(My-a)_+}+\norm{Ty-b}\bigr).
\end{equation}
The constant is uniform in the right-hand sides; other finite-dimensional
norms change only the constant, and equalities can be written as opposite
inequalities. The exact attainment proof needs only its existence for each
instance; the approximation proofs establish explicit encoding bounds. Neither argument gives a
uniform constant when the constraint normals vary with the leader.
